\documentclass[aspectratio=169,xcolor={dvipsnames,table}]{beamer}
\input{../../_en_preambulo_beamer.tex}
% Comandos y macros estadísticas compartidas para las presentaciones Beamer en inglés (Metropolis)

% Conjuntos numéricos básicos
\newcommand{\R}{\mathbb{R}}
\newcommand{\N}{\mathbb{N}}
\newcommand{\Q}{\mathbb{Q}}
\newcommand{\Z}{\mathbb{Z}}
\newcommand{\set}[1]{\left\{ #1 \right\}}
\newcommand{\abs}[1]{\left|#1\right|}
\newcommand{\norm}[1]{\left\| #1 \right\|}

% Comandos estadísticos y probabilísticos del curso
\newcommand{\Var}{\operatorname{Var}}
\newcommand{\cov}{\operatorname{Cov}}
\newcommand{\corr}{\rho}
\newcommand{\comb}[2]{\begin{pmatrix} #1 \\ #2 \end{pmatrix}}
\newcommand{\s}{\sigma}
\newcommand{\particion}{\mathfrak{P}}
\newcommand{\card}[1]{\#\left( #1 \right)}
\newcommand{\seq}[3]{\set{#1}_{#2}^{#3}}
\newcommand{\E}{\mathbb{E}}
\newcommand{\Prob}{\mathbb{P}}

% Entornos pedagógicos y cajas en inglés académico natural (soportando tanto nombres en inglés como en español)
\newenvironment{discussion}{%
  \begin{alertblock}{Discussion Question}%
}{%
  \end{alertblock}%
}
\newenvironment{discusion}{%
  \begin{alertblock}{Discussion Question}%
}{%
  \end{alertblock}%
}

\newenvironment{reflection}{%
  \begin{alertblock}{Classroom Reflection}%
}{%
  \end{alertblock}%
}
\newenvironment{reflexion}{%
  \begin{alertblock}{Classroom Reflection}%
}{%
  \end{alertblock}%
}

\newenvironment{paradox}{%
  \begin{alertblock}{Exploratory Paradox}%
}{%
  \end{alertblock}%
}
\newenvironment{paradoja}{%
  \begin{alertblock}{Exploratory Paradox}%
}{%
  \end{alertblock}%
}

\newenvironment{motivation}{%
  \begin{exampleblock}{Motivation: Data Science in Practice}%
}{%
  \end{exampleblock}%
}
\newenvironment{motivacion}{%
  \begin{exampleblock}{Motivation: Data Science in Practice}%
}{%
  \end{exampleblock}%
}

\newenvironment{quickchallenge}{%
  \begin{block}{Quick Team Challenge}%
}{%
  \end{block}%
}
\newenvironment{retoclase}{%
  \begin{block}{Quick Team Challenge}%
}{%
  \end{block}%
}


\title{Measures of Central Tendency}
\subtitle{Section 1.2 — Statistical Modeling for Data Science}
\author[J. Castillo Colmenares]{Juliho Castillo Colmenares}
\institute{Tecnológico de Monterrey}
\date{}

\begin{document}

% Slide 1: Title
\begin{frame}[plain]
  \titlepage
\end{frame}

% Slide 2: Chapter Roadmap and Position Indicator
\begin{frame}{Roadmap — Chapter 01: Descriptive Statistics}
  Where are we on our journey through foundational statistics? \pause
  \begin{itemize}
    \item \textbf{01.01 Introduction to Descriptive Statistics} \emph{(Completed)} \pause
    \item \color{TecRojo}\textbf{01.02 Measures of Central Tendency \emph{(Today)}}\color{black} \pause
    \item \textbf{01.03 Measures of Dispersion} \emph{(Next Session)}
  \end{itemize}
  \vspace{1em}
  \begin{block}{Section 1.2 Objective}
    Understand and calculate the center of a dataset mathematically (mean and pivot) and robustly in the presence of extreme noise (median and mode).
  \end{block}
\end{frame}

% Slide 3: Quick Team Challenge (Trigger Question)
\begin{frame}{Quick Team Challenge}
  \begin{alertblock}{Workplace Scenario (Trigger Question)}
    At a small artificial intelligence startup, four developers earn $\$30,000$ MXN per month each. Today, the company hires a new executive director with a monthly salary of $\$380,000$ MXN.
    \vspace{0.6em}
    \begin{enumerate}
      \item What is the average salary of all 5 company employees? \pause
      \item If human resources reports that average in job postings, are they giving an honest impression of what a typical employee earns? \pause
      \item What alternative descriptive measure would more honestly reflect the real office compensation?
    \end{enumerate}
  \end{alertblock}
\end{frame}

% Slide 4A: Classroom Reflection Conclusions (Analysis)
\begin{frame}{Classroom Reflection Conclusions — Analysis}
  \begin{itemize}
    \item \textbf{The Arithmetic Average:} Computation yields $\frac{4(30,000) + 380,000}{5} = \$100,000$ MXN per month. \pause
    \item \textbf{Outlier Distortion:} Not a single developer earns $\$100,000$. The executive's salary pulls the mean upward drastically, creating a statistical mirage. \pause
    \item \textbf{The Value of the Median:} Ordering salaries ($30, 30, \mathbf{30}, 30, 380$), the central value ($\$30,000$) represents the true employee experience. \pause
    \item \textbf{Key Lesson:} In Data Science, the mean alone is insufficient; we must know when to deploy robust measures.
  \end{itemize}
\end{frame}

% Slide 4B: Classroom Reflection Conclusions (Answers)
\begin{frame}{Classroom Reflection Conclusions — Answers}
  \begin{block}{Formal Answers to Quick Team Challenge}
    \begin{enumerate}
      \item \textbf{Arithmetic mean:} $\$100,000$ MXN per month. \pause
      \item \textbf{Is it representative?} No, it does not reflect the typical employee; $80\%$ of staff earn less than one-third of that figure. \pause
      \item \textbf{Honest metric:} The \textbf{Median ($\tilde{X}$)} or the \textbf{Mode ($M_o$)}, since they are immune to distortion by the star executive.
    \end{enumerate}
  \end{block}
\end{frame}

% Slide 5: Arithmetic Mean and Fundamental Property
\begin{frame}{The Arithmetic Mean ($\bar{X}$)}
  \scriptsize
  \vspace{-0.15cm}
  The \textbf{arithmetic mean} is the standard average of a set of $N$ observations $X_1, X_2, \dots, X_N$:
  \begin{align}
    \bar{X} = \dfrac{1}{N} \sum_{i=1}^{N} X_i
  \end{align}
  \pause
  \begin{block}{Center of Gravity Property}
    The algebraic sum of all deviations around the arithmetic mean is exactly equal to zero:
    \begin{align}
      \sum_{i=1}^{N} (X_i - \bar{X}) = 0
    \end{align}
    This proves that $\bar{X}$ acts as the exact physical balancing point (center of mass) of the distribution.
  \end{block}
\end{frame}

% Slide 6: The Pivot Concept (P)
\begin{frame}{The Pivot Concept ($P$)}
  \scriptsize
  \vspace{-0.15cm}
  When data values are large or clustered close together, we can simplify arithmetic operations using an arbitrary reference value called a \textbf{Pivot ($P$)}. \pause
  \begin{itemize}
    \item Let $d_i = X_i - P$ denote the deviation of each observation from the pivot $P$. \pause
    \item We can reconstruct the exact mean by adding the average of those deviations to the pivot:
  \end{itemize}
  \begin{align}
    \bar{X} = P + \dfrac{1}{N} \sum_{i=1}^{N} d_i
  \end{align}
  \pause
  \begin{block}{Algorithmic Advantage}
    This formula prevents floating-point overflow when summing millions of large numbers in computing, forming the basis for iterative mean algorithms in Big Data pipelines.
  \end{block}
\end{frame}

% Slide 7: Median and Mode
\begin{frame}{Robust Measures: Median and Mode}
  \scriptsize
  \vspace{-0.15cm}
  \begin{itemize}
    \item \textbf{The Median ($\tilde{X}$):} The value that divides the ordered sample into two exact halves ($50\%$ lower and $50\%$ upper).
      \begin{itemize}
        \item If $N$ is odd, it is the exact middle observation.
        \item If $N$ is even, it is the average of the two middle observations.
      \end{itemize}
    \pause
    \item \textbf{The Mode ($M_o$):} The value that appears most frequently in the dataset. It may not exist or there may be multiple modes (bimodal, multimodal). \pause
  \end{itemize}
  \vspace{0.4em}
  \begin{block}{Breakdown Point}
    The median features a $50\%$ breakdown point: up to half of the data can be corrupted or extreme outliers without distorting the median infinitely. The mean has a $0\%$ breakdown point.
  \end{block}
\end{frame}

% Slide 8: Central Tendency for Grouped Data
\begin{frame}{Central Tendency for Grouped Data}
  When processing frequency tables with $K$ classes or bins of relative/absolute frequencies: \pause
  \begin{itemize}
    \item Let $f_i$ be the absolute frequency of class $i$, and let $X_i$ be the \textbf{class mark} (interval midpoint). \pause
    \item The approximate mean for grouped data is weighted by each class frequency:
  \end{itemize}
  \begin{align}
    \bar{X}_{\text{grouped}} = \dfrac{\sum_{i=1}^{K} f_i X_i}{\sum_{i=1}^{K} f_i} = \dfrac{1}{n} \sum_{i=1}^{K} f_i X_i
  \end{align}
  \pause
  \begin{alertblock}{Practical Note}
    Grouped calculation is a highly accurate approximation to raw calculation, ideal when only aggregated histograms are preserved due to privacy or storage constraints.
  \end{alertblock}
\end{frame}

% Slide 9A: Live Python Lab (Classical Measures)
% Slide 9: Python Lab (Classical Measures)
\begin{frame}[fragile]{Python Lab: Classical Measures}
  \scriptsize
  Exploratory calculation of Arithmetic Mean, Median, and Mode using NumPy and SciPy.
  \\ Script: \texttt{\footnotesize presentations/code/01\_estadistica\_descriptiva/01.02\_mean\_median\_mode.py}
  
  \vspace{-0.1em}
  {\lstset{basicstyle=\ttfamily\tiny, aboveskip=0.2em, belowskip=0.2em}
  \lstinputlisting[firstline=1,lastline=16]{../../code/01_estadistica_descriptiva/01.02_mean_median_mode.py}
  }

  \vspace{-0.1em}
  \begin{itemize}
    \item Verify that the exact arithmetic mean (\texttt{np.mean}) matches the manual derivation.
    \item Observe that \texttt{stats.mode} returns the most frequent value along with its frequency count.
  \end{itemize}
\end{frame}

% Slide 11: Python Code — Computation via Pivot
\begin{frame}[fragile]{Code Laboratory: Computation via Pivot ($P = 9$)}
  \vspace{0.2em}
  {\lstset{basicstyle=\ttfamily\tiny, aboveskip=0.2em, belowskip=0.2em}
  \lstinputlisting[firstline=18,lastline=27]{../../code/01_estadistica_descriptiva/01.02_mean_median_mode.py}
  }

  \vspace{0.2em}
  \begin{block}{\footnotesize Verified Output (Console)}
    {\fontsize{6.5pt}{7.5pt}\selectfont
    \texttt{--- Classical Measures ---} \\
    \texttt{Arithmetic mean: 9.38 | Median: 9.50 | Mode: 5} \\
    \texttt{--- Pivot Calculation (P = 9) ---} \\
    \texttt{Deviations (d\_i): [-4 -1  2  0  3 -3  5  1] => Sum: 3} \\
    \texttt{Mean via pivot:   9 + (3/8) = 9.38}
    }
  \end{block}
\end{frame}

% Slide 10: Course Exercises - Level 1

% Slide 11: Course Exercises - Level 2

% Slide 12: Course Exercises - Level 3

% Slide 13: Course Exercises - Level 4

% Slide 14: Section 1.2 Summary
\begin{frame}{Section 1.2 Summary}
  \begin{block}{Mastered Concepts}
    \begin{itemize}
      \item \textbf{Numerical Center:} The arithmetic mean is the balance point ($\sum(X_i - \bar{X}) = 0$), but is vulnerable to extreme outliers. \pause
      \item \textbf{Pivot Efficiency:} $\bar{X} = P + \frac{\sum d_i}{N}$ simplifies manual calculation and stabilizes numerical computing algorithms. \pause
      \item \textbf{Statistical Robustness:} The median and mode isolate anomalies in skewed or outlier-heavy distributions.
    \end{itemize}
  \end{block}
  \pause
  \vspace{0.8em}
  \begin{center}
    \color{TecAzulOscuro}\textbf{Assimilation Pause:} Any questions regarding pivot calculation or choosing between mean and median?
    \color{black}
  \end{center}
\end{frame}

% Slide 15: Connection to Next Topic
\begin{frame}{Connection to the Next Topic}
  \begin{alertblock}{Next Section: 01.03 Measures of Dispersion}
    Two datasets can have exactly the same mean ($\bar{X} = 50$) yet behave completely differently ($[50, 50, 50]$ vs $[0, 50, 100]$).
    \vspace{0.6em}
    \\ In \textbf{Section 1.3}, we will measure how dispersed, variable, or spread out our data points are using:
    \begin{itemize}
      \item The Variance ($S^2$)
      \item The Standard Deviation ($S$)
      \item The Coefficient of Variation ($CV$)
    \end{itemize}
  \end{alertblock}
\end{frame}

\end{document}
